\chapter{Apéndices y Anexos}

%% \section{Apéndice A: Preguntas de la entrevista a cliente (Equipo de grupo MACS)}

\begin{enumerate}
    \item Desde sus experiencias con las plataformas USV del laboratorio, ¿cómo describirían su estado actual, y qué aspectos no se pueden modificar o deben preservarse al integrar hardware nuevo?
    \item ¿Qué tareas o misiones les gustaría que estas plataformas pudieran realizar que hoy no pueden?
    \item ¿Qué tipo de datos o resultados esperarían obtener al final de una misión?
    \item ¿Qué tipo de información ambiental les interesaría que el sistema pudiera medir durante una misión?
    \item ¿Dónde y cómo se instalarían físicamente los sensores y el hardware nuevo en el USV, y qué restricciones de peso, tamaño, consumo energético o pruebas de lastre debemos considerar?
    \item ¿Qué tan crítico es que los datos de los sensores estén sincronizados con la posición del vehículo en el momento de la captura?
    \item ¿Han considerado o descartado algún tipo de sensor de percepción hasta ahora? ¿Por qué?
    \item ¿Qué alcance y tipo de comunicación esperan entre los USVs y la estación base, considerando los cuerpos de agua de prueba?
    \item ¿Qué información necesitan recibir en tiempo real en la estación base durante una misión, y qué tan crítico es mantener esa conexión sin interrupciones?
    \item ¿Con qué frecuencia o continuidad necesitarían que se recolecten datos?
    \item Si tuvieran que definir cuándo una misión se considera “exitosa”, ¿qué tendría que pasar?
    \item ¿Qué margen de error en la ubicación o trayectoria, y qué porcentaje de fallas, considerarían aceptables antes de decir que el sistema no es confiable?
    \item ¿Existe algún estándar o referencia que ya usen para evaluar el desempeño de sistemas similares?
    \item ¿Qué condiciones del entorno, requisitos normativos o de seguridad, y permisos deben tenerse en cuenta para operar y probar el sistema, tanto en el laboratorio como en cuerpos de agua reales?
    \item Si un robot fallara o se perdiera durante una misión, ¿qué pasaría después? ¿existe alguna forma de localizarlo o recuperarlo?
    \item ¿Qué limitaciones de tiempo, presupuesto o recursos se deben considerar?
    \item Si tuvieran que elegir entre que el sistema sea más preciso o más fácil de operar, ¿cuál priorizarían, y por qué?
    \item Una vez finalizado este proyecto, ¿cómo se imaginan que el grupo seguiría usando o ampliando el sistema?
\end{enumerate}

